\documentclass[11pt,a4paper]{article}
\usepackage[utf8]{inputenc}
\usepackage[T1]{fontenc}
\usepackage[brazilian]{babel}
\usepackage{amsmath,amssymb,amsthm}
\usepackage{geometry}
\geometry{top=2.2cm,bottom=2.2cm,left=2cm,right=2cm}
\usepackage{booktabs}
\usepackage{array}
\usepackage{xcolor}
\usepackage{fancyhdr}
\usepackage{tikz}
\usetikzlibrary{trees,arrows.meta,positioning}
\usepackage{listings}
\usepackage{hyperref}
\hypersetup{
    colorlinks=true,
    linkcolor=blue!70!black,
    urlcolor=blue!70!black,
    citecolor=blue!70!black
}

\setlength{\headheight}{14.5pt}
\pagestyle{fancy}
\fancyhf{}
\lhead{\textbf{UFU -- FEELT} | Projeto e Análise de Algoritmos}
\rhead{Prof. Felipe A. Louza | \textbf{Lista 9}}
\lfoot{\small Arthur Pires Lopes -- 12221EEL001}
\rfoot{\thepage}

\setlength{\parindent}{0pt}
\setlength{\parskip}{5pt}

\lstset{
  basicstyle=\ttfamily\small,
  frame=single,
  rulecolor=\color{black!30},
  backgroundcolor=\color{black!3},
  numbers=left,
  numberstyle=\tiny\color{gray},
  stepnumber=1,
  numbersep=8pt,
  showstringspaces=false,
  breaklines=true,
  literate={á}{{\'a}}1 {ã}{{\~a}}1 {ç}{{\c c}}1 {é}{{\'e}}1 {í}{{\'i}}1 {ó}{{\'o}}1 {ú}{{\'u}}1 {Á}{{\'A}}1 {É}{{\'E}}1 {Í}{{\'I}}1 {Ó}{{\'O}}1 {Ú}{{\'U}}1 {≤}{{$\le$}}1 {≥}{{$\ge$}}1 {←}{{$\leftarrow$}}1 {≠}{{$\ne$}}1 {∞}{{$\infty$}}1
}

\begin{document}

% =========================================================================
% CABEÇALHO OFICIAL
% =========================================================================
\begin{center}
    {\Large \textbf{Universidade Federal de Uberlândia -- UFU}}\\[2pt]
    {\large Faculdade de Engenharia Elétrica (FEELT)}\\[2pt]
    {\large Projeto e Análise de Algoritmos (PAA) -- Prof. Felipe A. Louza}\\[4pt]
    \rule{\textwidth}{1pt}\\[5pt]
    {\LARGE \textbf{Lista 9: Programação Dinâmica}}\\[3pt]
    \textbf{Versão para Entrega Oficial}\\[2pt]
    \rule{\textwidth}{0.5pt}
\end{center}

\vspace{0.1cm}
\noindent \textbf{Aluno:} Arthur Pires Lopes \quad \textbf{Matrícula:} 12221EEL001

\vspace{0.3cm}

% =========================================================================
% QUESTÃO 1
% =========================================================================
\section*{Questão 1}
\textbf{Enunciado:} Mostre a execução passo-a-passo do algoritmo \textsc{Cut-Rod-Bottom-Up}$(P, n)$ para a seguinte entrada:
\begin{center}
\begin{tabular}{r|ccccccc}
\toprule
\textbf{tamanho} $i$ & 0 & 1 & 2 & 3 & 4 & 5 & 6 \\
\midrule
\textbf{preço} $p_i$ & 0 & 1 & 3 & 11 & 16 & 19 & 10 \\
\bottomrule
\end{tabular}
\end{center}

\subsection*{Resolução}

O algoritmo \textsc{Cut-Rod-Bottom-Up} calcula a receita máxima $A[i]$ que pode ser obtida cortando uma barra de comprimento $i \in \{0, 1, \dots, n\}$. A relação de recorrência baseada na subestrutura ótima é dada por:
\begin{equation}
A[0] = 0, \qquad A[i] = \max_{1 \le j \le i} \Big( p_j + A[i - j] \Big), \quad \text{para } 1 \le i \le n.
\end{equation}

Avaliamos iterativamente cada comprimento $i$ de $1$ até $n = 6$:

\begin{itemize}
    \item \textbf{Para $i = 0$:} Caso base, $A[0] = 0$.
    \item \textbf{Para $i = 1$:}
    \[
    j = 1: p_1 + A[0] = 1 + 0 = 1 \implies A[1] = 1.
    \]
    \item \textbf{Para $i = 2$:}
    \begin{align*}
    j = 1 &: p_1 + A[1] = 1 + 1 = 2 \\
    j = 2 &: p_2 + A[0] = 3 + 0 = 3 \implies A[2] = \max(2, 3) = 3.
    \end{align*}
    \item \textbf{Para $i = 3$:}
    \begin{align*}
    j = 1 &: p_1 + A[2] = 1 + 3 = 4 \\
    j = 2 &: p_2 + A[1] = 3 + 1 = 4 \\
    j = 3 &: p_3 + A[0] = 11 + 0 = 11 \implies A[3] = \max(4, 4, 11) = 11.
    \end{align*}
    \item \textbf{Para $i = 4$:}
    \begin{align*}
    j = 1 &: p_1 + A[3] = 1 + 11 = 12 \\
    j = 2 &: p_2 + A[2] = 3 + 3 = 6 \\
    j = 3 &: p_3 + A[1] = 11 + 1 = 12 \\
    j = 4 &: p_4 + A[0] = 16 + 0 = 16 \implies A[4] = \max(12, 6, 12, 16) = 16.
    \end{align*}
    \item \textbf{Para $i = 5$:}
    \begin{align*}
    j = 1 &: p_1 + A[4] = 1 + 16 = 17 \\
    j = 2 &: p_2 + A[3] = 3 + 11 = 14 \\
    j = 3 &: p_3 + A[2] = 11 + 3 = 14 \\
    j = 4 &: p_4 + A[1] = 16 + 1 = 17 \\
    j = 5 &: p_5 + A[0] = 19 + 0 = 19 \implies A[5] = \max(17, 14, 14, 17, 19) = 19.
    \end{align*}
    \item \textbf{Para $i = 6$:}
    \begin{align*}
    j = 1 &: p_1 + A[5] = 1 + 19 = 20 \\
    j = 2 &: p_2 + A[4] = 3 + 16 = 19 \\
    j = 3 &: p_3 + A[3] = 11 + 11 = \mathbf{22} \quad \text{(ótimo)} \\
    j = 4 &: p_4 + A[2] = 16 + 3 = 19 \\
    j = 5 &: p_5 + A[1] = 19 + 1 = 20 \\
    j = 6 &: p_6 + A[0] = 10 + 0 = 10 \implies A[6] = \max(20, 19, \mathbf{22}, 19, 20, 10) = 22.
    \end{align*}
\end{itemize}

\textbf{Tabela Final Preenchida:}
\begin{center}
\begin{tabular}{r|ccccccc}
\toprule
$i$ & 0 & 1 & 2 & 3 & 4 & 5 & 6 \\
\midrule
$p_i$ & 0 & 1 & 3 & 11 & 16 & 19 & 10 \\
$A[i]$ & \textbf{0} & \textbf{1} & \textbf{3} & \textbf{11} & \textbf{16} & \textbf{19} & \textbf{22} \\
\bottomrule
\end{tabular}
\end{center}

\textbf{Resultado:} O valor máximo obtido para a barra de tamanho $n = 6$ é $A[6] = 22$.

\vspace{0.4cm}

% =========================================================================
% QUESTÃO 2
% =========================================================================
\section*{Questão 2}
\textbf{Enunciado:} Escreva um algoritmo para descobrir os tamanhos dos pedaços nos quais a barra foi cortada na solução ótima. \emph{Dica: Modifique o algoritmo \textsc{Cut-Rod-Bottom-Up}$(P, n)$ para registrar os cortes na solução.}

\subsection*{Resolução}

Para reconstruir os pedaços da solução ótima, mantemos um vetor auxiliar $S[1 \dots n]$, onde $S[i]$ armazena o tamanho ótimo do primeiro pedaço $j$ cortado de uma barra de comprimento $i$.

\subsubsection*{1. Algoritmo Estendido com Registro de Cortes}

\begin{lstlisting}
EXTENDED-CUT-ROD-BOTTOM-UP(P, n)
1  let A[0 .. n] and S[1 .. n] be new arrays
2  A[0] = 0
3  for i = 1 to n
4      m = -infinity
5      for j = 1 to i
6          if m < P[j] + A[i - j]
7              m = P[j] + A[i - j]
8              S[i] = j
9      A[i] = m
10 return A and S
\end{lstlisting}

\subsubsection*{2. Algoritmo de Recuperação dos Pedaços}

\begin{lstlisting}
RETRIEVE-CUT-ROD-SOLUTION(P, n)
1  (A, S) = EXTENDED-CUT-ROD-BOTTOM-UP(P, n)
2  let L be an empty list
3  curr = n
4  while curr > 0
5      add S[curr] to L
6      curr = curr - S[curr]
7  return L
\end{lstlisting}

\subsubsection*{3. Rastreamento e Análise de Complexidade}
Aplicando à entrada da Questão 1 ($n = 6$):
\begin{center}
\begin{tabular}{r|cccccc}
\toprule
$i$ & 1 & 2 & 3 & 4 & 5 & 6 \\
\midrule
$S[i]$ & 1 & 2 & 3 & 4 & 5 & 3 \\
\bottomrule
\end{tabular}
\end{center}
Para $curr = 6$: $S[6] = 3 \implies$ adiciona 3, novo $curr = 6 - 3 = 3$. Para $curr = 3$: $S[3] = 3 \implies$ adiciona 3, novo $curr = 3 - 3 = 0$. Fim da execução. Pedaços ótimos: $\langle 3, 3 \rangle$ ($11 + 11 = 22$).
\begin{itemize}
    \item \textbf{Complexidade de Tempo:} $\Theta(n^2)$ para o preenchimento das tabelas e $O(n)$ no laço de recuperação. Total: $\Theta(n^2)$.
    \item \textbf{Complexidade de Espaço:} $\Theta(n)$ para os vetores auxiliares $A[0..n]$ e $S[1..n]$.
\end{itemize}

\vspace{0.4cm}

% =========================================================================
% QUESTÃO 3
% =========================================================================
\section*{Questão 3}
\textbf{Enunciado:} Mostre a execução passo-a-passo do algoritmo \textsc{LIS-PD}$(S, n)$ para a seguinte entrada:
\begin{center}
\begin{tabular}{r|cccccccc}
\toprule
$i$ & 1 & 2 & 3 & 4 & 5 & 6 & 7 & 8 \\
\midrule
$s_i$ & 5 & 2 & 8 & 6 & 3 & 6 & 9 & 7 \\
\bottomrule
\end{tabular}
\end{center}

\subsection*{Resolução}

O algoritmo calcula o vetor $L[1 \dots n]$, onde $L[i]$ denota o comprimento da maior subsequência estritamente crescente que \textbf{termina exatamente no elemento $s_i$}. A regra indutiva é:
\begin{equation}
L[1] = 1, \qquad L[i] = 1 + \max \left( \{ L[j] \mid 1 \le j < i \text{ e } s_j < s_i \} \cup \{0\} \right), \quad \text{para } 2 \le i \le n.
\end{equation}

Cálculo detalhado passo a passo:
\begin{itemize}
    \item \textbf{Para $i = 1$ ($s_1 = 5$):} $L[1] = 1$. Subsequência: $\langle 5 \rangle$.
    \item \textbf{Para $i = 2$ ($s_2 = 2$):}
    $j = 1: s_1 = 5 \not< 2$.
    Nenhum elemento anterior é menor $\implies \ell = 0 \implies L[2] = 1 + 0 = 1$. Subsequência: $\langle 2 \rangle$.
    \item \textbf{Para $i = 3$ ($s_3 = 8$):}
    \begin{align*}
    j = 1 &: s_1 = 5 < 8 \implies L[1] = 1 \\
    j = 2 &: s_2 = 2 < 8 \implies L[2] = 1 \implies \ell = \max(1, 1) = 1 \implies L[3] = 1 + 1 = 2.
    \end{align*}
    \item \textbf{Para $i = 4$ ($s_4 = 6$):}
    \begin{align*}
    j = 1 &: s_1 = 5 < 6 \implies L[1] = 1 \\
    j = 2 &: s_2 = 2 < 6 \implies L[2] = 1 \\
    j = 3 &: s_3 = 8 \not< 6 \implies \ell = 1 \implies L[4] = 1 + 1 = 2.
    \end{align*}
    \item \textbf{Para $i = 5$ ($s_5 = 3$):}
    \begin{align*}
    j = 1 &: s_1 = 5 \not< 3 \\
    j = 2 &: s_2 = 2 < 3 \implies L[2] = 1 \\
    j = 3, 4 &: s_3 = 8 \not< 3, \; s_4 = 6 \not< 3 \implies \ell = 1 \implies L[5] = 1 + 1 = 2.
    \end{align*}
    \item \textbf{Para $i = 6$ ($s_6 = 6$):}
    \begin{align*}
    j = 1 &: s_1 = 5 < 6 \implies L[1] = 1 \\
    j = 2 &: s_2 = 2 < 6 \implies L[2] = 1 \\
    j = 3, 4 &: s_3 = 8 \not< 6, \; s_4 = 6 \not< 6 \text{ (estritamente menor)} \\
    j = 5 &: s_5 = 3 < 6 \implies L[5] = 2 \implies \ell = \max(1, 2) = 2 \implies L[6] = 1 + 2 = 3.
    \end{align*}
    \item \textbf{Para $i = 7$ ($s_7 = 9$):}
    Candidatos $s_j < 9$: $s_1(1), s_2(1), s_3(2), s_4(2), s_5(2), s_6(3)$.
    \[
    \ell = \max(L[1], L[2], L[3], L[4], L[5], L[6]) = \max(1, 1, 2, 2, 2, 3) = 3 \implies L[7] = 1 + 3 = 4.
    \]
    \item \textbf{Para $i = 8$ ($s_8 = 7$):}
    Candidatos $s_j < 7$: $s_1=5(1), s_2=2(1), s_4=6(2), s_5=3(2), s_6=6(3)$. Note que $s_3 = 8 \not< 7$ e $s_7 = 9 \not< 7$.
    \[
    \ell = \max(1, 1, 2, 2, 3) = 3 \implies L[8] = 1 + 3 = 4.
    \]
\end{itemize}

\textbf{Tabela Final do Vetor $L$:}
\begin{center}
\begin{tabular}{r|cccccccc}
\toprule
$i$ & 1 & 2 & 3 & 4 & 5 & 6 & 7 & 8 \\
\midrule
$s_i$ & 5 & 2 & 8 & 6 & 3 & 6 & 9 & 7 \\
$L[i]$ & \textbf{1} & \textbf{1} & \textbf{2} & \textbf{2} & \textbf{2} & \textbf{3} & \textbf{4} & \textbf{4} \\
\bottomrule
\end{tabular}
\end{center}

\textbf{Resultado Final:} O tamanho máximo da LIS global é $\max_{1 \le j \le 8} L[j] = \mathbf{4}$ (atingido nas posições $i = 7$ e $i = 8$).

\vspace{0.4cm}

% =========================================================================
% QUESTÃO 4
% =========================================================================
\section*{Questão 4}
\textbf{Enunciado:} Escreva um algoritmo que retorna a LIS de uma sequência (não apenas o tamanho da LIS). \emph{Dica: Modifique o algoritmo \textsc{LIS-PD}$(S, n)$ para registrar as soluções parciais do problema.}

\subsection*{Resolução}

Adicionamos um vetor de ponteiros predecessores $P[1 \dots n]$. Sempre que encontramos um índice $j < i$ tal que $s_j < s_i$ que maximize $L[j]$, registramos $P[i] = j$. Ao término, identificamos o índice $k$ que maximiza $L[k]$ e realizamos o retrocesso (backtracking) encadeado.

\begin{lstlisting}
LIS-RECUPERA-SEQUENCIA(S, n)
1  let L[1 .. n] and P[1 .. n] be new arrays
2  L[1] = 1
3  P[1] = -1
4  max_len = 1
5  best_idx = 1
6  for i = 2 to n
7      l = 0
8      p = -1
9      for j = 1 to i - 1
10         if S[j] < S[i] and L[j] > l
11             l = L[j]
12             p = j
13     L[i] = 1 + l
14     P[i] = p
15     if L[i] > max_len
16         max_len = L[i]
17         best_idx = i
18 let R be an empty list
19 curr = best_idx
20 while curr != -1
21     add S[curr] to R
22     curr = P[curr]
23 reverse(R)
24 return R
\end{lstlisting}

\subsubsection*{Rastreamento e Análise}
Para a entrada da Questão 3 com $best\_idx = 7$ ($s_7 = 9$):
\[
P[7] = 6 \; (s_6 = 6) \to P[6] = 5 \; (s_5 = 3) \to P[5] = 2 \; (s_2 = 2) \to P[2] = -1.
\]
Invertendo a lista recuperada: $\langle 2, 3, 6, 9 \rangle$. (Alternativamente, para $best\_idx = 8$: $\langle 2, 3, 6, 7 \rangle$).
\begin{itemize}
    \item \textbf{Complexidade de Tempo:} $\sum_{i=1}^{n-1} i = \Theta(n^2)$ para preenchimento e $O(n)$ para reconstruir a lista. Total: $\Theta(n^2)$.
    \item \textbf{Complexidade de Espaço:} $\Theta(n)$ para armazenar $L[1..n]$ e $P[1..n]$.
\end{itemize}

\vspace{0.4cm}

% =========================================================================
% QUESTÃO 5
% =========================================================================
\section*{Questão 5}
\textbf{Enunciado:} Mostre a execução passo-a-passo do algoritmo \textsc{COINS-PD}$(C, k)$ para a seguinte entrada:
\[
C = [1, 5, 8, 11] \quad \text{e} \quad k = 13
\]

\subsection*{Resolução}

O algoritmo preenche um vetor $M[0 \dots k]$, onde $M[i]$ armazena o número mínimo de moedas para formar o valor $i$:
\begin{equation}
M[0] = 0, \qquad M[i] = 1 + \min_{\substack{1 \le j \le n \\ i - c_j \ge 0}} M[i - c_j], \quad \text{para } 1 \le i \le k.
\end{equation}

Execução iterativa para $i = 0 \dots 13$:
\begin{itemize}
    \item $i = 0$: $M[0] = 0$.
    \item $i = 1$: Moeda 1: $1 + M[0] = 1 \implies M[1] = 1$.
    \item $i = 2$: Moeda 1: $1 + M[1] = 2 \implies M[2] = 2$.
    \item $i = 3$: Moeda 1: $1 + M[2] = 3 \implies M[3] = 3$.
    \item $i = 4$: Moeda 1: $1 + M[3] = 4 \implies M[4] = 4$.
    \item $i = 5$: Moedas 1, 5: $1 + \min(M[4], M[0]) = 1 + \min(4, 0) = 1 \implies M[5] = 1$.
    \item $i = 6$: Moedas 1, 5: $1 + \min(M[5], M[1]) = 1 + \min(1, 1) = 2 \implies M[6] = 2$.
    \item $i = 7$: Moedas 1, 5: $1 + \min(M[6], M[2]) = 1 + \min(2, 2) = 3 \implies M[7] = 3$.
    \item $i = 8$: Moedas 1, 5, 8: $1 + \min(M[7], M[3], M[0]) = 1 + \min(3, 3, 0) = 1 \implies M[8] = 1$.
    \item $i = 9$: Moedas 1, 5, 8: $1 + \min(M[8], M[4], M[1]) = 1 + \min(1, 4, 1) = 2 \implies M[9] = 2$.
    \item $i = 10$: Moedas 1, 5, 8: $1 + \min(M[9], M[5], M[2]) = 1 + \min(2, 1, 2) = 2 \implies M[10] = 2$.
    \item $i = 11$: Moedas 1, 5, 8, 11: $1 + \min(M[10], M[6], M[3], M[0]) = 1 + \min(2, 2, 3, 0) = 1 \implies M[11] = 1$.
    \item $i = 12$: Moedas 1, 5, 8, 11: $1 + \min(M[11], M[7], M[4], M[1]) = 1 + \min(1, 3, 4, 1) = 2 \implies M[12] = 2$.
    \item $i = 13$: Moedas 1, 5, 8, 11:
    \[
    M[13] = 1 + \min\big( M[12], M[8], M[5], M[2] \big) = 1 + \min(2, 1, 1, 2) = 1 + 1 = \mathbf{2}.
    \]
\end{itemize}

\textbf{Tabela Final Preenchida:}
\begin{center}
\begin{tabular}{r|cccccccccccccc}
\toprule
$i$ & 0 & 1 & 2 & 3 & 4 & 5 & 6 & 7 & 8 & 9 & 10 & 11 & 12 & 13 \\
\midrule
$M[i]$ & \textbf{0} & \textbf{1} & \textbf{2} & \textbf{3} & \textbf{4} & \textbf{1} & \textbf{2} & \textbf{3} & \textbf{1} & \textbf{2} & \textbf{2} & \textbf{1} & \textbf{2} & \textbf{2} \\
\bottomrule
\end{tabular}
\end{center}

\textbf{Resultado:} O menor número de moedas para formar a quantia $k = 13$ é $M[13] = \mathbf{2}$ (composto por $8 + 5 = 13$).

\vspace{0.4cm}

% =========================================================================
% QUESTÃO 6
% =========================================================================
\section*{Questão 6}
\textbf{Enunciado:} Escreva um algoritmo que resolve o problema das moedas no qual \textbf{não é permitido repetições de moedas no troco}.

\subsection*{Resolução}

Quando cada moeda $c_i \in C$ pode ser utilizada \textbf{no máximo uma vez}, o problema torna-se análogo à Mochila 0-1. Definimos $M[i, v]$ como o número mínimo de moedas dentre o subconjunto $\{c_1, \dots, c_i\}$ necessárias para formar a quantia exata $v$.

\subsubsection*{1. Relação de Recorrência}
\begin{equation}
M[i, v] = 
\begin{cases}
0, & \text{se } v = 0 \\
\infty, & \text{se } i = 0 \text{ e } v > 0 \\
M[i-1, v], & \text{se } c_i > v \\
\min\Big(M[i-1, v], \; 1 + M[i-1, v - c_i]\Big), & \text{se } c_i \le v
\end{cases}
\end{equation}

\subsubsection*{2. Algoritmo com Otimização de Espaço ($1$D)}
Para economizar memória de $O(n \cdot k)$ para $O(k)$, percorremos o valor $v$ em \textbf{ordem decrescente} (de $k$ até $c_i$), impedindo que a mesma moeda seja reutilizada no mesmo estágio:

\begin{lstlisting}
COINS-01-KNAPSACK(C, n, k)
1  let dp[0 .. k] be a new array
2  dp[0] = 0
3  for v = 1 to k
4      dp[v] = infinity
5  for i = 1 to n
6      for v = k down to C[i]
7          if dp[v - C[i]] != infinity
8              dp[v] = min(dp[v], 1 + dp[v - C[i]])
9  if dp[k] == infinity
10     return -1    // Impossivel formar o troco sem repeticao
11 return dp[k]
\end{lstlisting}

\begin{itemize}
    \item \textbf{Complexidade de Tempo:} $\Theta(n \cdot k)$ (pseudo-polinomial).
    \item \textbf{Complexidade de Espaço:} $O(k)$ com vetor unidimensional (ou $\Theta(n \cdot k)$ se a matriz for preservada para recuperação das moedas).
\end{itemize}

\vspace{0.4cm}

% =========================================================================
% QUESTÃO 7
% =========================================================================
\section*{Questão 7}
\textbf{Enunciado:} Dados um conjunto de valores de moedas $C = c_1, c_2, \dots, c_n$ e uma quantia $k$. Escreva um algoritmo que retorne o \textbf{número de maneiras diferentes} que podemos formar a quantia $k$. \emph{Dica: Modifique o algoritmo \textsc{COINS-PD}$(C, k)$.}

\subsection*{Resolução}

Desejamos calcular o número de combinações não ordenadas de moedas cuja soma totaliza $k$ (onde moedas podem se repetir). 

\subsubsection*{1. Relação de Recorrência}
Seja $W[i, v]$ o número de maneiras de formar a quantia $v$ usando apenas moedas de $\{c_1, \dots, c_i\}$:
\begin{equation}
W[i, v] = W[i-1, v] + W[i, v - c_i], \quad \text{com } W[i, 0] = 1 \; (\forall i) \text{ e } W[0, v] = 0 \; (\forall v > 0).
\end{equation}

\subsubsection*{2. Algoritmo em Programação Dinâmica}
Mantendo o laço de moedas no nível externo e o laço de valores no nível interno crescente, evitamos permutações duplicadas (ex: $1+5$ e $5+1$ são contabilizados como uma única maneira):

\begin{lstlisting}
COUNT-WAYS-COIN-CHANGE(C, n, k)
1  let dp[0 .. k] be a new array
2  dp[0] = 1
3  for v = 1 to k
4      dp[v] = 0
5  for j = 1 to n
6      for v = C[j] to k
7          dp[v] = dp[v] + dp[v - C[j]]
8  return dp[k]
\end{lstlisting}

\begin{itemize}
    \item \textbf{Complexidade de Tempo:} $\Theta(n \cdot k)$ operações aritméticas.
    \item \textbf{Complexidade de Espaço:} $\Theta(k)$ no vetor acumulador.
\end{itemize}

\vspace{0.4cm}

% =========================================================================
% QUESTÃO 8
% =========================================================================
\section*{Questão 8}
\textbf{Enunciado:} Mostre a execução passo-a-passo do algoritmo \textsc{KNAPSACK-PD}$(n, K)$ para a seguinte entrada (capacidade $K = 10$):
\begin{center}
\begin{tabular}{c|cc}
\toprule
\textbf{item} $i$ & \textbf{peso} $w_i$ & \textbf{valor} $v_i$ \\
\midrule
1 & 6 & \$30 \\
2 & 3 & \$14 \\
3 & 4 & \$16 \\
4 & 2 & \$9 \\
\bottomrule
\end{tabular}
\end{center}

\subsection*{Resolução}

A matriz $M[1 \dots n, 0 \dots K]$ armazena o valor ótimo considerando os primeiros $i$ itens com capacidade $k$:
\begin{align}
M[1, k] &= \begin{cases} 0, & \text{se } w_1 > k \\ v_1, & \text{se } w_1 \le k \end{cases} \\
M[i, k] &= \begin{cases} M[i-1, k], & \text{se } k - w_i < 0 \\ \max\Big(M[i-1, k], \; M[i-1, k - w_i] + v_i\Big), & \text{se } k - w_i \ge 0 \end{cases}
\end{align}

\subsubsection*{Preenchimento Passo a Passo das Linhas}
\begin{itemize}
    \item \textbf{Linha 1 ($w_1 = 6, v_1 = 30$):}
    Para $k = 0 \dots 5$, $M[1, k] = 0$. Para $k = 6 \dots 10$, $M[1, k] = 30$.
    \item \textbf{Linha 2 ($w_2 = 3, v_2 = 14$):}
    \begin{itemize}
        \item $k = 0, 1, 2$: $k < 3 \implies M[2, k] = M[1, k] = 0$.
        \item $k = 3, 4, 5$: $\max(M[1, k], M[1, k-3] + 14) = \max(0, 0 + 14) = 14$.
        \item $k = 6, 7, 8$: $\max(M[1, k], M[1, k-3] + 14) = \max(30, 0 + 14) = 30$.
        \item $k = 9, 10$: $\max(M[1, k], M[1, k-3] + 14) = \max(30, 30 + 14) = 44$.
    \end{itemize}
    \item \textbf{Linha 3 ($w_3 = 4, v_3 = 16$):}
    \begin{itemize}
        \item $k = 0, 1, 2$: $M[3, k] = M[2, k] = 0$.
        \item $k = 3$: $k < 4 \implies M[3, 3] = M[2, 3] = 14$.
        \item $k = 4, 5$: $\max(M[2, k], M[2, k-4] + 16) = \max(14, 0 + 16) = 16$.
        \item $k = 6, 7, 8$: $\max(M[2, k], M[2, k-4] + 16) = \max(30, 0/14 + 16) = 30$.
        \item $k = 9$: $\max(M[2, 9], M[2, 5] + 16) = \max(44, 14 + 16 = 30) = 44$.
        \item $k = 10$: $\max(M[2, 10], M[2, 6] + 16) = \max(44, 30 + 16 = \mathbf{46}) = \mathbf{46}$.
    \end{itemize}
    \item \textbf{Linha 4 ($w_4 = 2, v_4 = 9$):}
    \begin{itemize}
        \item $k = 0, 1$: $0$.
        \item $k = 2$: $\max(0, 0 + 9) = 9$.
        \item $k = 3$: $\max(14, 0 + 9) = 14$.
        \item $k = 4$: $\max(16, 0 + 9) = 16$.
        \item $k = 5$: $\max(16, 14 + 9 = 23) = 23$.
        \item $k = 6, 7$: $\max(30, 16 + 9 = 25) = 30$.
        \item $k = 8$: $\max(30, 30 + 9 = 39) = 39$.
        \item $k = 9$: $\max(44, 30 + 9 = 39) = 44$.
        \item $k = 10$: $\max(46, 30 + 9 = 39) = \mathbf{46}$.
    \end{itemize}
\end{itemize}

\textbf{Matriz Completa $M[i, k]$:}
\begin{center}
\begin{tabular}{c|ccccccccccc}
\toprule
$i \setminus k$ & 0 & 1 & 2 & 3 & 4 & 5 & 6 & 7 & 8 & 9 & 10 \\
\midrule
1 ($w_1=6, v_1=30$) & 0 & 0 & 0 & 0 & 0 & 0 & \textbf{30} & 30 & 30 & 30 & 30 \\
2 ($w_2=3, v_2=14$) & 0 & 0 & 0 & 14 & 14 & 14 & 30 & 30 & 30 & 44 & 44 \\
3 ($w_3=4, v_3=16$) & 0 & 0 & 0 & 14 & 16 & 16 & 30 & 30 & 30 & 44 & \textbf{46} \\
4 ($w_4=2, v_4=9$)  & 0 & 0 & 9 & 14 & 16 & 23 & 30 & 30 & 39 & 44 & \textbf{46} \\
\bottomrule
\end{tabular}
\end{center}

\subsubsection*{Reconstrução da Solução Ótima por Retrocesso}
Partindo de $M[4, 10] = 46$:
\begin{enumerate}
    \item $M[4, 10] = 46 = M[3, 10] \implies$ \textbf{Item 4 NÃO entra}. Capacidade restante: $10$.
    \item $M[3, 10] = 46 \ne M[2, 10] = 44 \implies$ \textbf{Item 3 ENTRA}. Nova capacidade: $10 - w_3 = 10 - 4 = 6$.
    \item $M[2, 6] = 30 = M[1, 6] \implies$ \textbf{Item 2 NÃO entra}. Capacidade restante: $6$.
    \item $M[1, 6] = 30 > 0 \implies$ \textbf{Item 1 ENTRA}. Nova capacidade: $6 - w_1 = 6 - 6 = 0$.
\end{enumerate}

\textbf{Resultado:} Valor máximo = \textbf{\$46}, com os itens selecionados $\{1, 3\}$ e peso total $6 + 4 = 10\text{ kg}$.

\vspace{0.4cm}

% =========================================================================
% QUESTÃO 9
% =========================================================================
\section*{Questão 9}
\textbf{Enunciado:} Escreva um algoritmo que resolve o problema da mochila no qual \textbf{é permitido repetições de itens na mochila}.

\subsection*{Resolução}

No problema da \emph{Mochila Não Limitada} (Unbounded Knapsack), qualquer item $i \in \{1, \dots, n\}$ pode ser escolhido múltiplas vezes. 

\subsubsection*{1. Relação de Recorrência}
Seja $K[w]$ o valor máximo atingível com capacidade total $w \in \{0, \dots, W\}$:
\begin{equation}
K[0] = 0, \qquad K[w] = \max_{\substack{1 \le i \le n \\ w_i \le w}} \Big( K[w - w_i] + v_i \Big), \quad \text{para } 1 \le w \le W.
\end{equation}

\subsubsection*{2. Algoritmo em Programação Dinâmica}

\begin{lstlisting}
UNBOUNDED-KNAPSACK(w, v, n, W)
1  let K[0 .. W] be a new array
2  K[0] = 0
3  for c = 1 to W
4      K[c] = 0
5      for i = 1 to n
6          if w[i] <= c
7              K[c] = max(K[c], K[c - w[i]] + v[i])
8  return K[W]
\end{lstlisting}

\subsubsection*{3. Comparação Teórica com a Mochila 0-1}
Na formulação matricial bidimensional, a diferença crucial reside no índice da linha na transição:
\begin{itemize}
    \item \textbf{Mochila 0-1:} $M[i, w] = \max\Big(M[i-1, w], \; M[\mathbf{i-1}, w - w_i] + v_i\Big)$ (o item não pode mais ser usado).
    \item \textbf{Mochila com Repetições:} $M[i, w] = \max\Big(M[i-1, w], \; M[\mathbf{i}, w - w_i] + v_i\Big)$ (o mesmo item $i$ pode ser reutilizado).
\end{itemize}
\begin{itemize}
    \item \textbf{Complexidade de Tempo:} $\Theta(n \cdot W)$ (pseudo-polinomial).
    \item \textbf{Complexidade de Espaço:} $O(W)$ com o vetor unidimensional $K[0..W]$.
\end{itemize}

\vspace{0.4cm}

% =========================================================================
% QUESTÃO 10
% =========================================================================
\section*{Questão 10}
\textbf{Enunciado:} Escreva um algoritmo que retorna a LCS de duas sequências (não apenas o tamanho da LCS). \emph{Dica: Modifique o algoritmo \textsc{LCS-LENGTH}$(S, n)$ para registrar as soluções parciais do problema.}

\subsection*{Resolução}

Para retornar a subsequência comum de caracteres, modificamos o algoritmo para construir simultaneamente a matriz de tamanhos $M[0 \dots n, 0 \dots m]$ e uma matriz de setas/direções $B[1 \dots n, 1 \dots m]$, onde $B[i, j] \in \{\nwarrow, \uparrow, \leftarrow\}$.

\subsubsection*{1. Algoritmo Estendido com Registro de Rota}

\begin{lstlisting}
EXTENDED-LCS-LENGTH(X, n, Y, m)
1  let M[0 .. n, 0 .. m] and B[1 .. n, 1 .. m] be new matrices
2  for i = 0 to n
3      M[i, 0] = 0
4  for j = 1 to m
5      M[0, j] = 0
6  for i = 1 to n
7      for j = 1 to m
8          if X[i] == Y[j]
9              M[i, j] = M[i - 1, j - 1] + 1
10             B[i, j] = "DIAGONAL"  // Seta diagonal noroeste
11         else if M[i - 1, j] >= M[i, j - 1]
12             M[i, j] = M[i - 1, j]
13             B[i, j] = "UP"        // Seta para cima
14         else
15             M[i, j] = M[i, j - 1]
16             B[i, j] = "LEFT"      // Seta para a esquerda
17 return M and B
\end{lstlisting}

\subsubsection*{2. Algoritmo de Recuperação da Cadeia de Caracteres}

\begin{lstlisting}
PRINT-LCS(B, X, i, j)
1  if i == 0 or j == 0
2      return
3  if B[i, j] == "DIAGONAL"
4      PRINT-LCS(B, X, i - 1, j - 1)
5      print X[i]
6  else if B[i, j] == "UP"
7      PRINT-LCS(B, X, i - 1, j)
8  else
9      PRINT-LCS(B, X, i, j - 1)
\end{lstlisting}

\subsubsection*{3. Otimização de Espaço (Backtrack Direto sem Matriz $B$)}
Podemos prescindir integralmente da matriz $B$ inspecionando os valores diretamente em $M[i, j]$ durante o retrocesso:
\begin{lstlisting}
RETRIEVE-LCS-DIRECT(M, X, Y, n, m)
1  let L be an empty list
2  i = n, j = m
3  while i > 0 and j > 0
4      if X[i] == Y[j]
5          add X[i] to L
6          i = i - 1, j = j - 1
7      else if M[i - 1, j] >= M[i, j - 1]
8          i = i - 1
9      else
10         j = j - 1
11 reverse(L)
12 return L
\end{lstlisting}

\begin{itemize}
    \item \textbf{Complexidade de Tempo:} $\Theta(n \cdot m)$ para o preenchimento das matrizes e $O(n + m)$ para o retrocesso/reconstrução.
    \item \textbf{Complexidade de Espaço:} $\Theta(n \cdot m)$ para a matriz $M$ (economiza os $\Theta(n \cdot m)$ da matriz $B$ ao utilizar a reconstrução direta).
\end{itemize}

\end{document}
